%==============================================
% VOLUME I: BIT GENESIS
% (previously: "Substrate & Signal" / Part 1)
%==============================================
%
% NOTE ON VOLUME STRUCTURE:
%   Arliz is now organized as three volumes rather than three parts of
%   a single book (see Preface and Introduction). This file represents
%   the content of Volume I. If compiled as a standalone book, replace
%   \part{} below with \chapter{} of a suitable top-level structure, or
%   treat this file's content as the main matter of its own document.
%
% MERGE NOTES (step by step, voltage -> bits -> arrays):
%   - This volume keeps ALL original Part 1 chapters unchanged.
%   - No chapters from the old later parts are merged here: those
%     contain no purely "representation-level" content (everything
%     array-related from them belongs either to Volume II's hardware
%     material or to Volume III's array structures/algorithms/applications).
%   - The opening chapters ("The Nature of Data," "The Philosophy of
%     Representation") and "Voltage to Bits" remain the literal first
%     steps of the whole work's voltage -> bit -> number -> array progression.
%
\part{BIT GENESIS}
\label{part:bitgenesis}

\begin{partintro}
\lettrine[lines=3]{B}{efore we} can understand how arrays store elements, we must first understand how computers represent information itself. Every piece of data---numbers, text, images, instructions---exists as patterns of bits, and every bit begins its life as a voltage. This volume follows that single thread: from a voltage difference across a transistor, through binary encoding, to the numbers, characters, and bit-level structures that everything else in this work is built from.

\begin{quote}
\textit{``The choice of representation is often more important than the choice of algorithm.''}

\hfill--- \textsc{Donald E. Knuth}
\end{quote}
\end{partintro}

% --- Opening chapters: what data is, and what it means to represent it ---
\chapter{The Nature of Data}
What data actually is before it is encoded: the distinction between data, information, and meaning; raw measurements versus interpreted values; data as the record of a distinction or a difference; the role of an observer or interpreter in turning a physical trace into data; examples spanning natural "data" (tree rings, sediment layers, starlight) and human-made data (tally marks, ledgers, sensor readings); why data alone is inert without a scheme to read it; the seed of the idea that everything downstream---bits, numbers, arrays---is data wrapped in successive layers of agreed-upon meaning.

\chapter{The Philosophy of Representation}
What it means for one thing to stand for another; concrete versus abstract representation and the historical shift from tally marks to abstract number systems; Plato's allegory of the cave and the idea that representations can be incomplete or misleading; Aristotle's more practical view of mental representation as a useful internal model rather than a perfect copy; medieval and modern philosophical treatments of representation (Aquinas, Descartes, Kant, Sartre) and their relevance to encoding and interpretation; the layered nature of representation, where symbols represent symbols all the way down to physical reality; the abstraction hierarchy in computing, from physical voltages to applications, and where arrays sit within it; Shannon's information theory as the mathematical formalization of representation, encoding, decoding, and the bit as the fundamental unit of distinguishing information.

% --- Step 1: voltage and binary switching (the literal starting point) ---
\chapter{Voltage to Bits}
How transistors encode binary states, voltage levels, noise margins, signal integrity, why binary won over ternary.

\chapter{Binary Representation and Boolean Algebra}
Two-state logic, Boolean operations, truth tables, De Morgan's laws, bit as fundamental unit.

\chapter{Number Systems and Bases}
Positional notation, decimal, binary, octal, hexadecimal. Base conversion algorithms. Historical development. Why different bases matter for different purposes.

% --- Step 2: integers ---
\chapter{Integer Representation: Unsigned}
Binary positional notation, range limitations, conversion algorithms, hexadecimal convenience.

\chapter{Signed Integer Representation}
Sign-magnitude, one's complement, two's complement (why it won), arithmetic operations, overflow detection.

\chapter{Integer Overflow and Wraparound Behavior}
Undefined behavior, modular arithmetic, detecting overflow, saturating arithmetic, language-specific behaviors.

\chapter{Fixed-Point Representation}
Q-format notation, scaling factors, precision-range tradeoffs, embedded systems applications, fixed-point arithmetic.

% --- Step 3: floating point ---
\chapter{Floating-Point Representation: IEEE 754}
Sign, exponent, mantissa encoding, normalized and denormalized numbers, single vs. double precision.

\chapter{Special Floating-Point Values}
Infinity (positive and negative), NaN (quiet and signaling), signed zero, subnormal numbers.

\chapter{Floating-Point Arithmetic Operations}
Addition, multiplication, division, rounding modes (round to nearest, toward zero, toward), fused multiply-add.

\chapter{Floating-Point Error Analysis}
Precision loss, catastrophic cancellation, machine epsilon, relative error, Kahan summation algorithm.

\chapter{Decimal Floating-Point Representation}
IEEE 754-2008 decimal formats, financial computing requirements, densely packed decimal.

\chapter{Extended Precision and Arbitrary Precision}
80-bit extended precision, quadruple precision (128-bit), arbitrary precision libraries (GMP, MPFR).

% --- Step 4: characters and text ---
\chapter{Character Encoding: ASCII and Extensions}
7-bit ASCII, extended ASCII variants, code pages, limitations for internationalization.

\chapter{Unicode: Universal Character Encoding}
Code points, planes, combining characters, grapheme clusters, normalization forms (NFC, NFD, NFKC, NFKD).

\chapter{Unicode Transformation Formats}
UTF-8 (variable length, backward compatible), UTF-16 (surrogate pairs), UTF-32 (fixed length), BOM issues.

\chapter{Text Processing Complexities}
Grapheme vs. code point counting, case folding, locale-dependent operations, text segmentation.

% --- Step 5: byte ordering and exchange ---
\chapter{Endianness: Byte Ordering}
Big-endian vs. little-endian, historical reasons, network byte order, bi-endian systems, byte swapping.

\chapter{Cross-Platform Data Exchange}
Serialization formats, portable binary formats, protocol design considerations.

% --- Step 6: bit-level manipulation (directly feeds into Volume III bit arrays) ---
\chapter{Bitwise Operations Fundamentals}
AND, OR, XOR, NOT operations, truth tables, bit manipulation primitives.

\chapter{Bit Shifting and Rotation}
Logical shift, arithmetic shift, rotate left/right, applications to multiplication/division by powers of 2.

\chapter{Bit Manipulation Techniques}
Setting/clearing/toggling bits, bit masks, extracting bit fields, counting set bits (population count).

\chapter{Bit Packing and Flags}
Packing multiple boolean flags, bitfields in structs, union tricks, bit arrays.

\chapter{Advanced Bit Hacking}
Finding rightmost set bit, isolating bit patterns, bit reversal, Gray code conversion.

% --- Step 7: alignment (sets up cache-line discussion in Volume II) ---
\chapter{Data Alignment Fundamentals}
Natural alignment, alignment requirements by type, misalignment penalties.

\chapter{Structure Padding and Layout}
Compiler padding rules, struct member ordering, packing pragmas, cache line awareness.

\chapter{Alignment Optimization Techniques}
Manual padding, alignment attributes, reordering for cache efficiency.

% --- Step 8: media representations as arrays of encoded values ---
\chapter{Color Representation Models}
RGB, RGBA, HSV, HSL, CMYK, YUV, color spaces, gamma correction.

\chapter{Pixel Formats and Bit Depth}
1-bit, 8-bit indexed, 16-bit (RGB565), 24-bit, 32-bit (with alpha), HDR formats.

\chapter{Image Compression Fundamentals}
Lossless (PNG, GIF) vs. lossy (JPEG), compression ratios, quality tradeoffs.

\chapter{Video Encoding Principles}
Frame types (I, P, B), motion compensation, codecs (H.264, H.265, VP9, AV1).

\chapter{Audio Representation: Sampling Theory}
Nyquist theorem, sample rate, bit depth, quantization noise, aliasing.

\chapter{Audio Encoding Formats}
PCM (uncompressed), FLAC (lossless), MP3, AAC, Opus (lossy), perceptual coding.

\chapter{Signal Processing Representations}
Time domain, frequency domain, spectrograms, wavelet transforms.

% --- Step 9: pointers (the bridge to address calculation in Volume III) ---
\chapter{Pointer Representation}
How pointers are stored in memory, pointer size (32-bit vs. 64-bit), pointer tagging.

\chapter{Pointer Arithmetic and Address Calculation}
Array element addressing, struct member offsets, pointer subtraction.

\chapter{Special Pointer Values}
Null pointers, wild pointers, dangling pointers, pointer alignment requirements.

% --- Step 10: integrity and serialization (closing the representation layer) ---
\chapter{Error Detection Codes}
Parity bits (even/odd), checksums, CRC algorithms, hash-based integrity.

\chapter{Error Correction Codes}
Hamming codes, Reed-Solomon codes, LDPC codes, convolutional codes.

\chapter{Cyclic Redundancy Checks (CRC)}
CRC polynomials, CRC-8, CRC-16, CRC-32 standards, efficient implementation using tables, applications in data integrity.

\chapter{Data Serialization Fundamentals}
Converting in-memory structures to byte streams, portability issues.

\chapter{Binary Serialization Formats}
Protocol Buffers, FlatBuffers, Cap'n Proto, MessagePack, CBOR.

\chapter{Text Serialization Formats}
JSON, XML, YAML, TOML, human-readable vs. machine-efficient tradeoffs.

\chapter{Custom Binary Protocols}
Designing efficient binary formats, alignment considerations, versioning.

\chapter{Data Compression Theory}
Information theory basics, entropy, lossless vs. lossy bounds.

\chapter{Dictionary-Based Compression}
LZ77, LZ78, LZW, Lempel-Ziv family, dictionary construction.

\chapter{Entropy Coding}
Huffman coding, arithmetic coding, asymmetric numeral systems (ANS).

\chapter{Modern Compression Algorithms}
zlib, gzip, bzip2, LZMA, Zstandard, Brotli, compression levels.

\chapter{Specialized Compression}
Run-length encoding, delta encoding, columnar compression, domain-specific methods.